\documentclass[11pt]{article}
\usepackage[margin=1in]{geometry}
\usepackage{amsmath,amssymb,amsthm}
\usepackage{longtable,array,booktabs}
\usepackage{hyperref}
\newtheorem{theorem}{Theorem}
\newtheorem{lemma}[theorem]{Lemma}
\title{Detailed archival supplement:\\
A self-contained proof of primitive trace-one quadratics over \(\mathbb F_{2^k}\)}
\author{Baohua Li}
\date{2026-09-11}
\begin{document}
\maketitle
\begin{abstract}
This supplement is an archival, expanded version of the even-characteristic PTE proof used in Paper A. It contains the complete paper proof, additional derivations of the Gauss-sum and Ramanujan-sum ingredients, the full finite certificate for \(1\le k\le29\), the exact factorizations of \(2^{2k}-1\), and reproducibility and formalization notes. The theorem is known: the statement is the \(n=2,t=1\) case of Cohen (1990, Theorem~1, p.~1), while Moreno (1989) proved the trace-one quadratic case earlier. Cohen (2003, Theorem~1.5, p.~170) gives an equivalent square-order formulation. No priority is claimed.
\end{abstract}
\section{Status and scope}
The purpose of this note is to keep a detailed proof and computation archive for the following statement. It is not a claim of new priority for the theorem itself. Paper A uses a compressed version of this material in its Appendix~B.
\section{The theorem and notation}
\begin{theorem}\label{thm:pte}
For every \(k\ge1\), there exists
\[
\beta\in\mathbb F_{2^{2k}}^*,\qquad
\operatorname{ord}(\beta)=2^{2k}-1,\qquad
\operatorname{Tr}_{\mathbb F_{2^{2k}}/\mathbb F_{2^k}}(\beta)=1.
\]
Equivalently, there is an irreducible quadratic \(X^2+X+a\in\mathbb F_{2^k}[X]\) whose roots are primitive.
\end{theorem}
\section{Detailed proof}
\begin{proof}[Proof of Theorem~\ref{thm:pte}]
The statement is classical (see \cite{Mor} and
\cite[Theorem~1, p.~1]{Coh90}; it also follows from the equivalent square-order formulation in
\cite[Theorem~1.5, p.~170]{Coh}); we give a self-contained proof in the
even-characteristic case used above. This is a reconstruction of known results (Moreno; Cohen), and no priority is claimed; it is included only to make the paper self-contained.
Put $q=2^k$, $L=\mathbb{F}_{q^2}$, $K_0=\mathbb{F}_q$, and
$M=q^2-1=(q-1)(q+1)$. Let
\[
E_1=\{x\in L:x^q+x=1\}.
\]
The trace map $L\to K_0$ is surjective and $K_0$-linear, so $|E_1|=q$. If
$x\in E_1$, then $x\notin K_0$ and its minimal polynomial over $K_0$ is
$X^2+X+xx^q$. Hence it is enough to prove that $E_1$ contains an element of
order $M$.

Let $\psi$ be the nontrivial additive character of $\mathbb F_2$, and put
$\psi_L(x)=\psi(\operatorname{Tr}_{L/\mathbb F_2}(x))$ and
$\psi_{K_0}(a)=\psi(\operatorname{Tr}_{K_0/\mathbb F_2}(a))$. For a multiplicative
character $\chi$ of $L^*$ define
\[
g_L(\chi)=\sum_{x\in L^*}\chi(x)\psi_L(x),\qquad
g_{K_0}(\chi)=\sum_{a\in K_0^*}\chi(a)\psi_{K_0}(a),\qquad
S(\chi)=\sum_{x\in E_1}\chi(x).
\]
If $\chi$ is nontrivial on $L^*$, then $|g_L(\chi)|=q$; if
$\chi|_{K_0^*}$ is nontrivial, then $|g_{K_0}(\chi)|=\sqrt q$. Indeed, the usual
squared-modulus calculation gives
\[
|g_L(\chi)|^2
=\sum_{t\in L^*}\chi(t)^{-1}\sum_{a\in L^*}\psi_L(a(1-t))=q^2,
\]
and the same argument in $K_0$ gives $q$.

We first compute $S(\chi)$. For $\alpha\in K_0$ set
$E_\alpha=\{x\in L:x^q+x=\alpha\}$; these sets partition $L$, and
$E_\alpha=\alpha E_1$ for $\alpha\ne0$. If $\chi^{q+1}\ne1$, then the
restriction of $\chi$ to $K_0^*$ is nontrivial; hence
\[
g_L(\chi)
=\sum_{\alpha\in K_0}\psi_{K_0}(\alpha)\sum_{x\in E_\alpha}\chi(x)
=S(\chi)g_{K_0}(\chi),
\]
so $|S(\chi)|=q/\sqrt q=\sqrt q$.

Suppose now that $\chi^{q+1}=1$ and $\chi\ne1$. Then $\chi|_{K_0^*}=1$.
Let $H=\{h\in L^*:h^{q+1}=1\}$; since $\gcd(q-1,q+1)=1$, the group $L^*$
is the direct product $K_0^*\times H$. Write $x=ch$ with $c\in K_0^*$ and
$h\in H$. The condition $x^q+x=1$ becomes
$c(h+h^{-1})=1$; thus $h\ne1$ and $c=(h+h^{-1})^{-1}$. This gives a
bijection between $H\setminus\{1\}$ and $E_1$, and therefore
\[
S(\chi)=\sum_{h\in H,\ h\ne1}\chi(h)
=\sum_{h\in H}\chi(h)-1=-1,
\]
because $\chi|_H$ is nontrivial.

We shall use the primitive-element indicator
\begin{equation}\label{eq:pte-indicator}
\mathbf 1_{\operatorname{ord}(x)=M}
=\frac{\varphi(M)}M
\sum_{\delta\mid M}\frac{\mu(\delta)}{\varphi(\delta)}
\sum_{\operatorname{ord}(\chi)=\delta}\chi(x).
\end{equation}
To see this, fix a generator $g$ of $L^*$ and write $x=g^a$. The left side
is the indicator of $\gcd(a,M)=1$. Its Fourier coefficient at a character
of order $\delta$ is the Ramanujan sum
$c_M(j)=\mu(\delta)\varphi(M)/\varphi(\delta)$; Fourier inversion gives
\eqref{eq:pte-indicator}. (The displayed identity for $c_M(j)$ follows by
M\"obius inversion.)

Assume for contradiction that $E_1$ has no element of order $M$. Summing
\eqref{eq:pte-indicator} over $E_1$ gives
\[
0=\frac{\varphi(M)}M
\sum_{\delta\mid M}\frac{\mu(\delta)}{\varphi(\delta)}S_\delta,
\qquad
S_\delta=\sum_{\operatorname{ord}(\chi)=\delta}S(\chi).
\]
Let $s=\omega(M)$ and $s'=\omega(q+1)$ be the numbers of distinct prime
divisors. The $\delta=1$ term is $q$. If $\delta\mid q+1$ and $\delta>1$,
then each character of order $\delta$ satisfies $\chi^{q+1}=1$, so the
corresponding contribution is $-\mu(\delta)$; summing over
$\delta\mid q+1$, $\delta>1$, gives $1$. For the remaining squarefree
divisors $\delta\nmid q+1$, the character sums are bounded by $\sqrt q$,
and their number is at most $2^s-2^{s'}$. Consequently,
\[
q+1\le(2^s-2^{s'})\sqrt q. \tag{A.2}
\]
Since $q+1>\sqrt q$, this implies
\[
2^s>2^{k/2}, \qquad\text{so}\qquad s>k/2.
\]
In particular $k<2s$. Let $p_1<p_2<\cdots$ be the primes and
$P_s=p_1\cdots p_s$. Since $M=2^{2k}-1$ has exactly $s$ distinct prime
divisors, $M\ge P_s$. We have
\[
P_{16}=32589158477190044730>2^{64}.
\]
For $n\ge16$, the ratio
\[
\frac{P_{n+1}/2^{4(n+1)}}{P_n/2^{4n}}
=\frac{p_{n+1}}{16}
\]
is greater than $1$, because $p_{n+1}\ge59$. Hence $P_s>2^{4s}$ for
every $s\ge16$. But $k<2s$ gives
\[
M=2^{2k}-1<2^{2k}<2^{4s},
\]
so $s\ge16$ is impossible. Thus $s\le15$ and $k<2s\le30$, i.e.
$k\le29$.

It remains to cover this finite range. In the table below, $m_k$ is an
irreducible polynomial of degree $k$ over $\mathbb F_2$, encoded by the
binary coefficient vector of the displayed hexadecimal integer; thus
$K_0=\mathbb F_2[x]/(m_k)$. The element $a_k\in K_0$ is encoded similarly.
For each row let $t$ be a root of $X^2+X+a_k$ in
$K_0[t]/(X^2+X+a_k)$. The deterministic verifier checks that $m_k$ is
irreducible, that $\operatorname{Tr}_{K_0/\mathbb F_2}(a_k)=1$, that
$t^{2^k}=t+1$, and that $t^M=1$ while $t^{M/p}\ne1$ for every prime
$p\mid M$. The last two checks show that $\operatorname{ord}(t)=M$; the
third shows that $\operatorname{Tr}_{L/K_0}(t)=1$. A verifier for the table, \texttt{verify\_k1\_29.py},
is included as an ancillary file; its SHA-256 hash is
\texttt{042CE7F74EDFF13ADE5C6285E42CB3D15491FC25}\allowbreak
\texttt{C65B99BC670D256047EE8D9F}.
\begin{center}
\small
\begin{tabular}{rrr}
$k$ & $m_k$ & $a_k$\\
\hline
1 & \texttt{0x3} & \texttt{0x1}\\
2 & \texttt{0x7} & \texttt{0x3}\\
3 & \texttt{0xb} & \texttt{0x5}\\
4 & \texttt{0x13} & \texttt{0xd}\\
5 & \texttt{0x25} & \texttt{0x13}\\
6 & \texttt{0x43} & \texttt{0x34}\\
7 & \texttt{0x83} & \texttt{0x47}\\
8 & \texttt{0x11b} & \texttt{0xc8}\\
9 & \texttt{0x203} & \texttt{0x2b}\\
10 & \texttt{0x409} & \texttt{0xcb}\\
11 & \texttt{0x805} & \texttt{0x525}\\
12 & \texttt{0x1009} & \texttt{0x39e}\\
13 & \texttt{0x201b} & \texttt{0x1ec4}\\
14 & \texttt{0x4021} & \texttt{0x3717}\\
15 & \texttt{0x8003} & \texttt{0x62bf}\\
16 & \texttt{0x1002b} & \texttt{0x1e6e}\\
17 & \texttt{0x20009} & \texttt{0x1155b}\\
18 & \texttt{0x40009} & \texttt{0x2ee84}\\
19 & \texttt{0x80027} & \texttt{0xb9af}\\
20 & \texttt{0x100009} & \texttt{0x6aee2}\\
21 & \texttt{0x200005} & \texttt{0x15750b}\\
22 & \texttt{0x400003} & \texttt{0x24a150}\\
23 & \texttt{0x800021} & \texttt{0x6c92f9}\\
24 & \texttt{0x100001b} & \texttt{0x389ec4}\\
25 & \texttt{0x2000009} & \texttt{0x6801ad}\\
26 & \texttt{0x400001b} & \texttt{0x2707d2e}\\
27 & \texttt{0x8000027} & \texttt{0x6e9ab1e}\\
28 & \texttt{0x10000003} & \texttt{0xfac7cd2}\\
29 & \texttt{0x20000005} & \texttt{0xc69b55e}\\
\end{tabular}
\end{center}
The table therefore supplies the required element for every $k\le29$.
Combining this with the contradiction for $k\ge30$ proves the theorem.
\end{proof}

\section{Supplementary derivations}
\subsection{Gauss-sum squared modulus}
For a nontrivial multiplicative character \(\chi\) of \(L^*\), put
\[
g_L(\chi)=\sum_{x\in L^*}\chi(x)\psi_L(x).
\]
Then
\[
|g_L(\chi)|^2
=\sum_{t\in L^*}\chi(t)^{-1}\sum_{a\in L^*}\psi_L(a(1-t)).
\]
For \(t=1\), the inner sum is \(|L|-1\); for \(t\ne1\), it is \(-1\). Since \(\sum_{t\in L^*}\chi(t)^{-1}=0\), the double sum is \(|L|\). This is the calculation used in the proof.
\subsection{Ramanujan-sum identity}
Let \(M\) be the order of \(L^*\), fix a generator \(g\), and write \(x=g^a\). For a character \(\chi_j(g)=e^{2\pi i j/M}\), the Fourier coefficient of the indicator of \(\gcd(a,M)=1\) is
\[
c_M(j)=\sum_{\substack{a\bmod M\\(a,M)=1}}e^{-2\pi ija/M}.
\]
With \(\delta=M/\gcd(M,j)\), the identity
\[
c_M(j)=\sum_{d\mid\gcd(M,j)}d\,\mu(M/d)=\mu(\delta)\frac{\varphi(M)}{\varphi(\delta)}
\]
follows by Möbius inversion. Fourier inversion then gives the primitive-element indicator used in the paper. This is the detailed form of the formula cited in the main proof.
\subsection{Why the finite checks suffice}
For each row of the certificate, let \(t\) be a root of \(X^2+X+a_k\) in \(K_0[t]/(X^2+X+a_k)\). The verifier checks:
\begin{enumerate}
\item \(m_k\) is irreducible of degree \(k\), so \(K_0=\mathbb F_2[x]/(m_k)\cong\mathbb F_{2^k}\);
\item \(\operatorname{Tr}_{K_0/\mathbb F_2}(a_k)=1\), so \(X^2+X+a_k\) is irreducible over \(K_0\);
\item \(t^{2^k}=t+1\), hence \(\operatorname{Tr}_{L/K_0}(t)=1\);
\item \(t^M=1\) and \(t^{M/p}\ne1\) for every prime \(p\mid M\).
\end{enumerate}
The last two conditions imply \(\operatorname{ord}(t)=M=2^{2k}-1\). Thus \(t\) is the required primitive trace-one element.
\section{Full finite certificate}
The following table gives the hexadecimal polynomial encodings used by the verifier:
\[
m_k=\sum_{i=0}^k m_{k,i}x^i,\qquad a_k=\sum_{i=0}^{k-1}a_{k,i}x^i,
\]
with the hexadecimal integer encoding the binary coefficient vector.

\begin{longtable}{rrrr}
\toprule
\(k\) & \(m_k\) & \(a_k\) & \(\omega(2^{2k}-1)\)\\
\midrule
\endfirsthead
\toprule
\(k\) & \(m_k\) & \(a_k\) & \(\omega(2^{2k}-1)\)\\
\midrule
\endhead
1 & \texttt{0x3} & \texttt{0x1} & 1 \\
2 & \texttt{0x7} & \texttt{0x3} & 2 \\
3 & \texttt{0xB} & \texttt{0x5} & 2 \\
4 & \texttt{0x13} & \texttt{0xD} & 3 \\
5 & \texttt{0x25} & \texttt{0x13} & 3 \\
6 & \texttt{0x43} & \texttt{0x34} & 4 \\
7 & \texttt{0x83} & \texttt{0x47} & 3 \\
8 & \texttt{0x11B} & \texttt{0xC8} & 4 \\
9 & \texttt{0x203} & \texttt{0x2B} & 4 \\
10 & \texttt{0x409} & \texttt{0xCB} & 5 \\
11 & \texttt{0x805} & \texttt{0x525} & 4 \\
12 & \texttt{0x1009} & \texttt{0x39E} & 6 \\
13 & \texttt{0x201B} & \texttt{0x1EC4} & 3 \\
14 & \texttt{0x4021} & \texttt{0x3717} & 6 \\
15 & \texttt{0x8003} & \texttt{0x62BF} & 6 \\
16 & \texttt{0x1002B} & \texttt{0x1E6E} & 5 \\
17 & \texttt{0x20009} & \texttt{0x1155B} & 3 \\
18 & \texttt{0x40009} & \texttt{0x2EE84} & 8 \\
19 & \texttt{0x80027} & \texttt{0xB9AF} & 3 \\
20 & \texttt{0x100009} & \texttt{0x6AEE2} & 7 \\
21 & \texttt{0x200005} & \texttt{0x15750B} & 6 \\
22 & \texttt{0x400003} & \texttt{0x24A150} & 7 \\
23 & \texttt{0x800021} & \texttt{0x6C92F9} & 4 \\
24 & \texttt{0x100001B} & \texttt{0x389EC4} & 9 \\
25 & \texttt{0x2000009} & \texttt{0x6801AD} & 7 \\
26 & \texttt{0x400001B} & \texttt{0x2707D2E} & 7 \\
27 & \texttt{0x8000027} & \texttt{0x6E9AB1E} & 6 \\
28 & \texttt{0x10000003} & \texttt{0xFAC7CD2} & 8 \\
29 & \texttt{0x20000005} & \texttt{0xC69B55E} & 6 \\
\bottomrule
\end{longtable}

The exact prime-power factorizations used by the verifier are:

\begin{longtable}{r p{1.2cm} p{11cm}}
\toprule
\(k\) & \(\omega(M)\) & factorization of \(M=2^{2k}-1\)\\
\midrule
\endfirsthead
\toprule
\(k\) & \(\omega(M)\) & factorization of \(M=2^{2k}-1\)\\
\midrule
\endhead
1 & 1 & \(3\) \\
2 & 2 & \(3\cdot 5\) \\
3 & 2 & \(3^{2}\cdot 7\) \\
4 & 3 & \(3\cdot 5\cdot 17\) \\
5 & 3 & \(3\cdot 11\cdot 31\) \\
6 & 4 & \(3^{2}\cdot 5\cdot 7\cdot 13\) \\
7 & 3 & \(3\cdot 43\cdot 127\) \\
8 & 4 & \(3\cdot 5\cdot 17\cdot 257\) \\
9 & 4 & \(3^{3}\cdot 7\cdot 19\cdot 73\) \\
10 & 5 & \(3\cdot 5^{2}\cdot 11\cdot 31\cdot 41\) \\
11 & 4 & \(3\cdot 23\cdot 89\cdot 683\) \\
12 & 6 & \(3^{2}\cdot 5\cdot 7\cdot 13\cdot 17\cdot 241\) \\
13 & 3 & \(3\cdot 2731\cdot 8191\) \\
14 & 6 & \(3\cdot 5\cdot 29\cdot 43\cdot 113\cdot 127\) \\
15 & 6 & \(3^{2}\cdot 7\cdot 11\cdot 31\cdot 151\cdot 331\) \\
16 & 5 & \(3\cdot 5\cdot 17\cdot 257\cdot 65537\) \\
17 & 3 & \(3\cdot 43691\cdot 131071\) \\
18 & 8 & \(3^{3}\cdot 5\cdot 7\cdot 13\cdot 19\cdot 37\cdot 73\cdot 109\) \\
19 & 3 & \(3\cdot 174763\cdot 524287\) \\
20 & 7 & \(3\cdot 5^{2}\cdot 11\cdot 17\cdot 31\cdot 41\cdot 61681\) \\
21 & 6 & \(3^{2}\cdot 7^{2}\cdot 43\cdot 127\cdot 337\cdot 5419\) \\
22 & 7 & \(3\cdot 5\cdot 23\cdot 89\cdot 397\cdot 683\cdot 2113\) \\
23 & 4 & \(3\cdot 47\cdot 178481\cdot 2796203\) \\
24 & 9 & \(3^{2}\cdot 5\cdot 7\cdot 13\cdot 17\cdot 97\cdot 241\cdot 257\cdot 673\) \\
25 & 7 & \(3\cdot 11\cdot 31\cdot 251\cdot 601\cdot 1801\cdot 4051\) \\
26 & 7 & \(3\cdot 5\cdot 53\cdot 157\cdot 1613\cdot 2731\cdot 8191\) \\
27 & 6 & \(3^{4}\cdot 7\cdot 19\cdot 73\cdot 87211\cdot 262657\) \\
28 & 8 & \(3\cdot 5\cdot 17\cdot 29\cdot 43\cdot 113\cdot 127\cdot 15790321\) \\
29 & 6 & \(3\cdot 59\cdot 233\cdot 1103\cdot 2089\cdot 3033169\) \\
\bottomrule
\end{longtable}

\section{Reproducibility}
The deterministic verifier is \texttt{verify\_k1\_29.py}, stored in the Paper D \texttt{verification/} directory and also included in the submission package. It uses only integer and polynomial arithmetic and requires neither network access nor SymPy. Run:
\begin{verbatim}
python verification/verify_k1_29.py
\end{verbatim}
The expected final line is:
\begin{verbatim}
All certificates verified.
\end{verbatim}
The SHA-256 of the verifier is
\begin{center}
\texttt{042CE7F74EDFF13ADE5C6285E42CB3D15491FC25C65B99BC670D256047EE8D9F}
\end{center}
The verifier was also independently reimplemented during the red-team audit; all \(k=1,\dots,29\) rows passed irreducibility, trace, Frobenius, and exact-order checks.
\section{Formalization boundary}
The affine group-generation argument and selected algebraic lemmas are formalized in Lean~4 with mathlib. A conditional Lean file formalizes the field-level Frobenius identities and the order-theoretic \(\tau\)-collapse once an explicit order interface is supplied. The analytic PTE proof and the finite \(k\le29\) certificate in this note are not Lean-formalized. This is a disclosed boundary, not a hidden dependency.
\section{Historical attribution}
The theorem is classical. The exact historical theorem used here is the \(n=2,t=1\) case of \cite[Theorem~1, p.~1]{Coh90}. Moreno \cite{Mor} proved the trace-one quadratic case earlier. Cohen \cite[Theorem~1.5, p.~170]{Coh} gives an equivalent square-order formulation. These references are historical attributions only; the proof above is self-contained.
\begin{thebibliography}{9}
\bibitem{Coh} S.~D.~Cohen, \emph{The orders of related elements of a finite field}, Ramanujan J.~7 (2003) 169--183, doi:10.1023/A:1026295128600.
\bibitem{Coh90} S.~D.~Cohen, \emph{Primitive elements and polynomials with arbitrary trace}, Discrete Math.~83 (1990) 1--7.
\bibitem{Mor} O.~Moreno, \emph{On the existence of a primitive quadratic of trace \(1\) over \(GF(p^m)\)}, J.~Combin.~Theory Ser.~A 51 (1989) 104--110.
\end{thebibliography}
\end{document}
